\documentclass[11pt,a4paper]{article}

% Compileer met pdfLaTeX. Installeer in MiKTeX zo nodig tcolorbox.
\usepackage[utf8]{inputenc}
\usepackage[T1]{fontenc}
\usepackage{lmodern}
\usepackage{xcolor}
\usepackage[breakable]{tcolorbox}

\definecolor{studieblauw}{HTML}{2563A6}
\setlength{\parindent}{0pt}
\setlength{\parskip}{5pt}
\setlength{\textwidth}{160mm}
\setlength{\oddsidemargin}{0mm}
\setlength{\evensidemargin}{0mm}

% Per begrip: korte betekenis, gevolgd door punten die je moet kennen.
\newtcolorbox{definition}[1]{
  title=\textbf{#1},
  colback=studieblauw!4,
  colframe=studieblauw,
  boxrule=0.7pt,
  arc=1.5mm,
  breakable,
  before skip=9pt,
  after skip=9pt
}

\newcommand{\kennen}{\par\smallskip\textbf{Kennen:}}

\title{Business Processes\\Introductie -- begrippenlijst}
\author{Dries Van den Brande}
\date{}

\begin{document}
\maketitle

\noindent\textit{Studielijst bij ``0. BP Intro''. Leer de betekenis en de
verschillen tussen begrippen; probeer de vragen onderaan zonder de lijst te
beantwoorden.}

\section{Bedrijfsprocessen en IT}

\begin{definition}{Bedrijfsproces}
Een reeks samenhangende activiteiten, gericht op een klant en afgestemd op de
organisatiedoelen.
\kennen
\begin{itemize}
  \item Herken een proces als een samenhang van activiteiten, niet als een losse taak.
  \item Benoem wie de klant is en welke waarde het proces voor die klant oplevert.
  \item Leg bij een toepassing uit welk bedrijfsproces ze ondersteunt.
\end{itemize}
\end{definition}

\begin{definition}{IT in een bedrijfsproces}
Applicaties ondersteunen bedrijfsprocessen. IT kan ook zelf deel uitmaken van
het aangeboden product of de dienst.
\kennen
\begin{itemize}
  \item Kijk verder dan de code: wie gebruikt de toepassing en voor welk werk?
  \item Herken beide rollen van IT in een concreet voorbeeld.
\end{itemize}
\end{definition}

\section{Drie types bedrijfsprocessen}

\begin{definition}{Primair bedrijfsproces}
Een operationeel proces dat rechtstreeks verband houdt met het doel van de
organisatie en rechtstreeks waarde cre\"eert voor de externe klant.
\kennen
\begin{itemize}
  \item Bepaal wat de organisatie aan haar klanten levert.
  \item Herken de primaire activiteit in een concreet bedrijf.
\end{itemize}
\end{definition}

\begin{definition}{Besturend bedrijfsproces}
Een proces dat de strategie en koers van de organisatie bepaalt (richten) en de
organisatie op haar doelen afstemt (inrichten).
\kennen
\begin{itemize}
  \item Onderscheid \textbf{richten} van \textbf{inrichten}.
  \item Leg uit waarom managementactiviteiten besturend kunnen zijn.
\end{itemize}
\end{definition}

\begin{definition}{Ondersteunend bedrijfsproces}
Een proces dat het algemene functioneren van de organisatie mogelijk maakt.
\kennen
\begin{itemize}
  \item Herken voorbeelden zoals beveiliging, catering, HR, ICT en finance.
  \item Bepaal wie binnen de organisatie de ondersteuning gebruikt.
\end{itemize}
\end{definition}

\begin{definition}{Externe klant}
Een klant buiten de organisatie die waarde uit haar aanbod ontvangt.
\kennen
\begin{itemize}
  \item Herken de externe klant van het primaire proces.
  \item Voorbeeld: de klant die een mobiele app laat ontwikkelen.
\end{itemize}
\end{definition}

\begin{definition}{Interne klant}
Een medewerker of afdeling binnen de organisatie die een proces gebruikt om
het eigen werk mogelijk te maken.
\kennen
\begin{itemize}
  \item Herken interne klanten van besturende en ondersteunende processen.
  \item Voorbeeld: collega's die HR- of ICT-ondersteuning nodig hebben.
\end{itemize}
\end{definition}

\section{Soorten organisaties en waardecreatie}

\begin{definition}{Productgerichte organisatie}
Een organisatie die fysieke of digitale producten maakt, verhandelt of
verdeelt.
\kennen
\begin{itemize}
  \item Onderscheid productie, handel en distributie.
  \item Bepaal per soort wat de primaire activiteit is.
\end{itemize}
\end{definition}

\begin{definition}{Productiebedrijf}
Een organisatie die producten vervaardigt. Productie is daarbij een primair
proces; waarde ontstaat door transformatie.
\kennen
\begin{itemize}
  \item Voorbeeld: een producent die tomaten tot tomatensaus verwerkt.
  \item Leg uit wat er door het productieproces verandert.
\end{itemize}
\end{definition}

\begin{definition}{Handelsbedrijf}
Een organisatie die producten aankoopt en verkoopt.
\kennen
\begin{itemize}
  \item Herken aankoop, opslag en verkoop als primaire processen.
  \item Leg uit hoe bewaren, verzamelen en distribueren waarde toevoegen.
\end{itemize}
\end{definition}

\begin{definition}{Distributiebedrijf}
Een organisatie die producten verdeelt of vervoert binnen de keten van
producent tot consument.
\kennen
\begin{itemize}
  \item Herken het verschil tussen een product maken, verhandelen en vervoeren.
  \item Voorbeeld uit de slides: een voedseltransportbedrijf.
\end{itemize}
\end{definition}

\begin{definition}{Bedrijfskolom}
Het geheel van organisaties dat producten van grondstof tot consument brengt.
\kennen
\begin{itemize}
  \item Zet de schakels van een eenvoudig product in de juiste volgorde.
  \item Gebruik de voorbeelden gepelde garnalen en katoenen T-shirt om de keten te oefenen.
\end{itemize}
\end{definition}

\begin{definition}{Dienstverlenende organisatie}
Een organisatie die een dienst, expertise of ervaring aan klanten levert.
\kennen
\begin{itemize}
  \item Onderscheid het leveren van een dienst van het leveren van een product.
  \item Herken dienstverlening aan organisaties (B2B) en consumenten (B2C).
\end{itemize}
\end{definition}

\begin{definition}{Kenmerken van een dienst}
Een dienst is niet tastbaar, vergankelijk, interactief en heterogeen.
\kennen
\begin{itemize}
  \item \textbf{Niet tastbaar:} je kunt de dienst niet vastpakken; na de ervaring is ze voorbij.
  \item \textbf{Vergankelijk:} je kunt een dienst niet bewaren.
  \item \textbf{Interactief:} aanbieder en klant werken samen.
  \item \textbf{Heterogeen:} de situatie verschilt van keer tot keer.
\end{itemize}
\end{definition}

\begin{definition}{Overheid en semi-publieke organisaties}
Organisaties die waarde leveren zonder winstoogmerk als centraal doel.
\kennen
\begin{itemize}
  \item Herken de voorbeelden uit de slides: gemeenten, Vlaamse en federale overheid,
        ziekenhuizen, vzw's en ngo's.
  \item Bedenk wie de ontvanger van hun waarde is.
\end{itemize}
\end{definition}

\section{Evolutie: van bezit naar toegang}

\begin{definition}{As a service}
Een model waarbij een product steeds vaker als dienst wordt aangeboden: de
klant krijgt toegang of gebruik in plaats van alleen bezit.
\kennen
\begin{itemize}
  \item Herken een product dat je vandaag als dienst gebruikt.
  \item Geef de redenen uit de slides voor deze verschuiving:
        wederkerende inkomsten, milieuvriendelijker, een beheersbare uitgave
        voor de klant en gemakkelijker te distribueren.
  \item Benoem welk bedrijfsproces achter het gekozen voorbeeld zit.
\end{itemize}
\end{definition}

\section{Toepassen zonder spiekbrief}

\begin{enumerate}
  \item Wat maakt een reeks activiteiten tot een bedrijfsproces?
  \item Wat is het verschil tussen primair, besturend en ondersteunend?
  \item Een IT-agency ontwikkelt mobiele apps en heeft developers, management,
        HR, sales en finance. Welke activiteiten horen bij welk procestype?
  \item Wie is de externe klant van die IT-agency? Noem ook twee interne klanten.
  \item Wat is een bedrijfskolom? Teken er een voor een katoenen T-shirt.
  \item Welk verschil zie je tussen de waardecreatie van een productiebedrijf en
        een handelsbedrijf?
  \item Noem de vier kenmerken van een dienst en illustreer er twee.
  \item Welk product gebruik jij ``as a service'', en welk proces maakt dat
        aanbod mogelijk?
\end{enumerate}

\end{document}
